% Generated draft: requires mathematical and source review.
\documentclass[11pt]{article}
\usepackage[margin=25mm]{geometry}
\usepackage{fontspec}
\setmainfont{Latin Modern Roman}
\usepackage{amsmath,amssymb,mathtools,unicode-math}
\setmathfont{Latin Modern Math}
\usepackage{enumitem,needspace,xurl,hyperref}
\setlength{\parindent}{0pt}
\setlength{\parskip}{5pt}
\setlength{\emergencystretch}{3em}
\newcommand{\N}{\mathbb{N}}
\newcommand{\Z}{\mathbb{Z}}
\newcommand{\Q}{\mathbb{Q}}
\newcommand{\R}{\mathbb{R}}
\newcommand{\C}{\mathbb{C}}
\newcommand{\F}{\mathbb{F}}
\newcommand{\A}{\mathbb{A}}
\newcommand{\PP}{\mathbb{P}}
\newcommand{\E}{\mathbb{E}}
\newcommand{\eps}{\varepsilon}
\newcommand{\dd}{\,\mathrm{d}}
\begin{document}
\section*{UnsolvedMath ID: 20002715}
\textbf{Source problem code:} catalogue entry
\subsection{Definitions and mathematical statement}
Charney's Problem 4.1 has three classification questions: (1) which arbitrary posets admit a Garside labeling; (2) how inequivalent labelings of one Garside poset are related, and when the labeling is unique; (3) when a given edge-labeled poset embeds in a Garside structure, when the corresponding monoids and groups embed, and what the minimal obstructions are. The source does not fix one precise labeling/embedding convention or equivalence relation.

The aid treats a finite faithful divisor-interval convention. A Garside monoid is atomic and left- and right-cancellative, admits left and right least common multiples and greatest common divisors, and has a balanced element $\Delta$ whose finite common set of left/right divisors generates it. A faithful labeling here identifies the poset with this divisor interval, labels a cover by its right quotient atom, and uses the associated path-relation monoid. Equivalence permits both a poset automorphism and a bijective renaming of labels. This explicit convention is not asserted to encompass every possible meaning of the source's terms.

Let $D_m$ be the height-two lattice with minimum $0$, maximum $1$ and $m\ge1$ incomparable atoms $a_1,\ldots,a_m$. Label $0\lessdot a_i$ by distinct $x_i$. The reported classification says the top labels are exactly
\[
 \lambda(a_i\lessdot1)=x_{\sigma(i)}\qquad(\sigma\in S_m),
\]
with associated monoid and group presentations
\[
 M_\sigma=\langle x_1,\ldots,x_m\mid
           x_i x_{\sigma(i)}=x_jx_{\sigma(j)}\text{ for all }i,j\rangle^+,
\]
\[
 G_\sigma=\langle x_1,\ldots,x_m\mid
           x_i x_{\sigma(i)}=x_jx_{\sigma(j)}\text{ for all }i,j\rangle.
\]
The common length-two element is $\Delta$. Under this convention the aid reports that $M_\sigma$ is Garside, has the prescribed divisor interval and embeds in its group of fractions $G_\sigma$. Two labelings are equivalent exactly when their permutations are conjugate. Hence $D_m$ has $p(m)$ equivalence classes, where $p(m)$ is the number of integer partitions of $m$.

Finite chains have a unique faithful labeling under the same convention: for a chain of $r$ covers, one atom labels every cover and $\Delta=x^r$. Thus the four-element diamond $D_2$ is the smallest finite bounded lattice with nonunique faithful labeling. For its transposition permutation, $G_{(12)}=\langle x,y\mid xy=yx\rangle\cong\Z^2$; for its identity permutation, $G_{\mathrm{id}}=\langle x,y\mid x^2=y^2\rangle$ is nonabelian. These distinct groups show why unlabeled-poset shape alone does not determine the associated group.

The result classifies only this height-two family and finite chains. It gives no general decision criterion for arbitrary posets, no classification of all inequivalent Garside structures and no complete set of embedding obstructions. In particular the count $p(m)$ depends on allowing poset automorphisms in equivalence; a pointwise-fixed-poset convention is a different question.
\subsection{Short English statement}
Classify Garside-labelable posets, their labelings and embedding obstructions. The aid solves a specific faithful height-two case: labelings correspond to permutation conjugacy classes, and the four-element diamond already supports inequivalent structures with different groups.
\subsection{Sources}
\begin{itemize}
\item R. Charney, Problem 4.1, American Institute of Mathematics workshop \emph{Braid groups, clusters and free probability}, page 8 as identified by supplied verification, \url{https://aimath.org/WWN/braidgroups/braidgroups.pdf}.
\item Local UnsolvedMath registry, permanent ID 20002715, source code \texttt{catalogue entry}, canonical location \texttt{aim-probability-notes.json}, \texttt{notes[156]}, record dated 14 August 2026. The corrected-verified statement repairs the split number 4.1, lost spaces and line-break hyphenation only; the original literature field is empty.
\item Supplied research attempt 1: partial result, problem status at run open, rigor labelled proved, novelty confidence low. Its faithful-labeling classification is machine-generated research aid, not peer reviewed; independently verify before citation. No separate authoritative reference for this special case is supplied.
\end{itemize}
\end{document}
